% ECONOMICS_PROBLEM: {"id": "H21-405", "title": "Optimal top-income taxation with avoidance and mobility", "jel_code": "H21", "source_id": "OP-0041"}
% FIRST_POSTED: 2026-10-09
% ATTEMPT_STATUS: partial
% ENTRY_KIND: research_attempt
\documentclass[11pt]{article}
\usepackage[margin=1in]{geometry}
\usepackage[utf8]{inputenc}
\usepackage{amsmath,amssymb,amsthm,hyperref}
\newtheorem{theorem}{Theorem}
\newtheorem{proposition}[theorem]{Proposition}
\newtheorem{lemma}[theorem]{Lemma}
\newtheorem{corollary}[theorem]{Corollary}
\theoremstyle{definition}
\newtheorem{example}[theorem]{Example}
\newtheorem{remark}[theorem]{Remark}
\newcommand{\E}{\mathbb E}
\newcommand{\R}{\mathbb R}
\title{Optimal top-income taxation with avoidance and mobility}
\author{GPT 6 Astra Ultra}
\date{9 October 2026}
\begin{document}
\maketitle
\section*{Research target and a restricted accounting result}
The target optimizes a top marginal rate jointly with enforcement and
reclassification policy, accounting for all displaced tax bases and a fixed
initial population. A local elasticity of reported top income cannot identify
that global optimum. We derive a transparent two-base marginal formula, show
why identical reported-income responses imply different welfare effects, and
provide a conditional robust-policy certificate. The formula is a benchmark,
not a universal tax recommendation.

\section*{A two-base economy with real and shifting responses}
Use one unit mass of identical taxpayers for the first calculation. Let $L$
be real earnings above a normalized zero threshold, $S$ earnings reclassified
to a second base taxed at fixed rate $\sigma$, and $X=L-S$ the top-rate base.
For fixed enforcement $e$, each taxpayer chooses feasible $(L,S)$ to maximize
\[
 V(\tau,e)=\max_{L,S}\{(1-\tau)L+(\tau-\sigma)S-c(L)-k(S,e)\}. \tag{1}
\]
Assume an interior unique smooth optimum, with convex real effort cost $c$
and shifting resource cost $k$. Taxpayer utility is quasilinear in the common
numeraire. The top rate changes no other primitive. Revenue and welfare are
\[
 R=\tau X+\sigma S,\qquad W=gV+\lambda[R-C(e)],
\]
where $g\ge0$ is a fixed social weight and $\lambda>0$ is the public-funds weight.
All receipts are included once. In particular $V$ is already after-tax private
welfare; adding pretax income separately would double-count it.

\begin{proposition}[Marginal welfare with cross-base recapture]
At a differentiability point, holding $e$ and $\sigma$ fixed,
\[
 \frac{\partial W}{\partial\tau}
 =(\lambda-g)X+\lambda\{\tau L_\tau-(\tau-\sigma)S_\tau\}. \tag{2}
\]
\end{proposition}
\begin{proof}
The envelope theorem applied to (1) gives $V_\tau=-L+S=-X$.
Differentiate total revenue to obtain
$R_\tau=X+\tau(L_\tau-S_\tau)+\sigma S_\tau$.
Substitution proves (2). The effort and shifting-resource effects already enter
$V$ through optimization, so they must not be subtracted again.
\end{proof}
Real earnings reductions destroy revenue at marginal rate $\tau$, whereas
pure shifting loses only the rate difference $\tau-\sigma$. This distinction
is an accounting consequence conditional on the model, not an empirical claim
that shifting is harmless. The resource cost $k$ and public enforcement cost
$C$ remain present in welfare.

\section*{Same observed elasticity, different marginal welfare}
At a common baseline let $X=100$, $\tau=0.5$, $\sigma=0.4$ and
$X_\tau=-100$. Consider local model A with $L_\tau=-99,S_\tau=1$ and model B
with $L_\tau=-1,S_\tau=99$. The reported tax base and its derivative coincide.
With $\lambda=1,g=0.7$, (2) equals $30-49.5-0.1=-19.6$ in A and $30-0.5-9.9=19.6$ in B.
These local responses can be generated by the quadratic primitives in the next
section, choosing intercepts to match baseline $X$. Thus even the sign of the
marginal welfare gain is not identified by the reported-income elasticity.

A mobility margin adds a further distinction. If the resident mass is $N(\tau)$,
per-resident bases are $x,s$, and revenue is
$N(\tau)(\tau x+\sigma s)$, then
\[
 R_\tau=N\{x+\tau x_\tau+\sigma s_\tau\}
       +N_\tau(\tau x+\sigma s).                         \tag{3}
\]
The last term is lost domestic revenue, while the welfare of emigrants stays
inside a fixed-initial-population criterion. Dropping them from the utility
integral changes the question. Their optimized migration costs and foreign
incomes require a residence model; neither (2) nor a taxable-income elasticity
identifies them. We do not add (3) to (2) while silently retaining its fixed-mass
envelope formula.

\section*{A computable joint rate/enforcement benchmark}
Let
\[
 c(L)=\frac{(L-L_0)^2}{2b},\quad
 k(S,e)=\frac{(\kappa_0+e)S^2}{2},\quad
 C(e)=\gamma e^2/2,
\]
with $b,\kappa_0,\gamma>0$. Restrict the policy to a compact rectangle
$\tau\in[\sigma,\bar\tau]$, $e\in[0,\bar e]$, $\bar\tau<1$,
and choose constants ensuring $0\le S<L$ on it. Then strict concavity of (1)
gives
\[
 L=L_0+b(1-\tau),\qquad S=\frac{\tau-\sigma}{\kappa_0+e}.
\]
The optimized utility is
\[
 V=(1-\tau)L_0+\frac b2(1-\tau)^2+
               \frac{(\tau-\sigma)^2}{2(\kappa_0+e)}.
\]
Substitute these rational expressions into $W$ to obtain a continuous function
on the rectangle, so an optimum exists. From the envelope and revenue formulas,
\[
 W_e=-gS^2/2+\lambda(\tau-\sigma)\frac{S}{\kappa_0+e}
                    -\lambda\gamma e.                   \tag{4}
\]
The first term records the extra real private shifting cost for given choices;
the second records revenue recovered because shifting falls. Its sign is not
fixed by saying enforcement reduces avoidance.

There is no need to assume global concavity to certify an approximate optimum.
Bound the two absolute partial derivatives of the explicit $W$ over the compact
rectangle by $L_\tau,L_e$, using $\kappa_0+e\ge\kappa_0$. On a rectangular
mesh with spacings $h_\tau,h_e$, choose the largest computed value, with each
value error at most $\eta$. Rounding any policy to a mesh point changes welfare
by at most $L_\tau h_\tau+L_e h_e$. Hence the chosen policy has regret at most
\[
 L_\tau h_\tau+L_e h_e+2\eta.                             \tag{5}
\]
This is a finite global certificate, albeit computationally expensive in a
higher-dimensional institution-specific problem.

\section*{Identified models and robust choice}
If parameters lie in a data-consistent compact set $\Theta$, define
$\underline W(p)=\inf_{\theta\in\Theta}W(p;\theta)$ and
$\overline W(p)=\sup_{\theta\in\Theta}W(p;\theta)$.
For any candidate policy $p_*$,
\[
 \sup_{\theta\in\Theta}
 [\sup_pW(p;\theta)-W(p_*;\theta)]
 \le\sup_p\overline W(p)-\underline W(p_*).              \tag{6}
\]
Proof: each term in brackets is bounded by the right side before taking the
supremum. This sufficient certificate is conservative because its upper and
lower endpoints may use different models. If every $W$ is within $\delta$ of
a fitted $\widehat W$ uniformly and $p_*$ is $\eta$-optimal for the latter,
its regret is at most $2\delta+\eta$ by adding and subtracting fitted values.
Neither bound removes uncertainty without credible data restrictions.

Linked all-base records, independently varied enforcement, migration histories
and a consistent welfare boundary are necessary to evaluate the expanded
problem. Bargaining rent transfers require modeling the recipients who lose
those rents. The present proof addresses real responses and reclassification
but not that channel. The substantive unresolved task is to justify a joint
identified class and obtain useful small regret under those richer margins.

\section{Completion Estimate}
45\%

This is a post hoc, heuristic estimate for the full catalogue target.
The attempt proves a two-base marginal welfare formula, demonstrates elasticity-based welfare ambiguity, and certifies joint rate and enforcement optimization in a quadratic benchmark. A data-supported joint model covering mobility, bargaining transfers, evasion, and all displaced tax bases remains necessary for the full policy target.

\begin{thebibliography}{9}
\bibitem{source} T. Piketty, E. Saez and S. Stantcheva (2014),
\emph{Optimal Taxation of Top Labor Incomes: A Tale of Three Elasticities},
AEJ: Economic Policy 6(1), 230--271.
\url{https://www.aeaweb.org/articles?id=10.1257/pol.6.1.230}.
Publisher abstract verifies the real, avoidance and bargaining distinctions;
(1)--(6) are a separately specified and derived benchmark, not the paper's
full optimal-tax model or a claimed new general formula.
\end{thebibliography}

\end{document}
